% Auto-generated by src/analysis/latex_tables.py — do not edit by hand.
\begin{table}[htbp]
\centering
\small
\caption{Pretreatment-characteristic time-trend robustness (Bai 2009; Hsiao 2014)}
\label{tab:pretreatment_trends}
\begin{tabular}{lcccc}
\toprule
 & CBR & Marriage rate & $I_g$ & $N$ \\
 & (1) & (2) & (3) & \\
\midrule
(1) Baseline (no pretreatment trends) & -0.0040 & -0.0042$^{***}$ & -0.00000 & 10,777 \\
 & (0.0032) & (0.0009) & (0.00006) & \\
(2) + literacy x trend & -0.0053 & -0.0043$^{***}$ & -0.00003 & 9,638 \\
 & (0.0034) & (0.0010) & (0.00006) & \\
(3) + lit + urban x trend & -0.0054 & -0.0043$^{***}$ & -0.00003 & 9,638 \\
 & (0.0034) & (0.0010) & (0.00006) & \\
(4) + lit + urban + pruss x trend & -0.0054 & -0.0043$^{***}$ & -0.00003 & 9,638 \\
 & (0.0034) & (0.0010) & (0.00006) & \\
(5) + lit + urban + pruss + jew x trend & 0.0002 & -0.0022$^{**}$ & 0.00005 & 9,638 \\
 & (0.0035) & (0.0009) & (0.00006) & \\
(6) + lit + urban + pruss + jew + women-15-49-share x trend & -0.0001 & -0.0022$^{**}$ & 0.00005 & 9,638 \\
 & (0.0034) & (0.0009) & (0.00006) & \\
(7) + born-in-Kreis share x trend & 0.0010 & -0.0023$^{***}$ & 0.00005 & 9,638 \\
 & (0.0034) & (0.0009) & (0.00006) & \\
(8) + 1849 student share x trend & 0.0021 & -0.0015 & 0.00009 & 7,143 \\
 & (0.0039) & (0.0010) & (0.00007) & \\
(9) + land Gini 1882 + log income-tax-pc 1876 x trend & 0.0002 & -0.0021$^{**}$ & 0.00006 & 7,143 \\
 & (0.0040) & (0.0011) & (0.00007) & \\
\bottomrule
\end{tabular}
\begin{minipage}{\linewidth}
\vspace{0.5em}
\footnotesize \textit{Notes:} Each row is a separate two-way fixed-effects regression of the outcome on $\mathrm{CathShare} \times \mathrm{Post}$ and the listed baseline iPEHD-1871 characteristics interacted with a centred linear time trend. The marital-fertility column reports the recalibrated Coale $I_g$ (Hutterite-normalised legitimate births per married woman 15--49, with the STA1871 marriage-prevalence shifter); see \S6.5 of the data appendix. The interactions allow counties with different pre-treatment literacy ($\mathrm{school1517}$), urbanisation ($f_{\mathrm{urban}}$), Prussian-citizenship share ($f_{\mathrm{pruss}}$), and Jewish-population share ($f_{\mathrm{jew}}$) to follow different trajectories. The Kulturkampf coefficient is then identified from deviations from those trajectories at 1873. The marriage-rate coefficient is essentially unchanged when literacy, urbanisation, and Prussian citizenship are added (rows 2--4) but attenuates by ~50\% when Jewish share is added (row 5), which absorbs differential dynamics in eastern provinces with high Jewish settlement and thus partially captures the same Polish-province channel documented in Tables~\ref{tab:falsifications} and \ref{tab:emigration_robustness}. Sample shrinks slightly because the iPEHD merge covers $\sim$90\% of Galloway counties. Standard errors clustered at the county level. $^{*}\,p<0.10$, $^{**}\,p<0.05$, $^{***}\,p<0.01$.
\end{minipage}
\end{table}
